\chapter{Stand zu Projektbeginn}

\section{Ausgangszustand des Projekts}

Das im Rahmen dieser Studienarbeit weiterentwickelte Robotersystem basiert auf einer zuvor durchgeführten Studienarbeit, deren Ziel die Konzeption und der Aufbau eines funktionsfähigen Demonstrators für den Einsatz auf Hochschulmessen war. Im Mittelpunkt dieser ersten Projektphase standen die Entwicklung eines Gesamtkonzepts, die Auswahl geeigneter Hard- und Softwarekomponenten sowie die Validierung grundlegender Funktionen des Roboters. \cite{...}

Im Rahmen der vorangegangenen Arbeit wurde zunächst eine Anforderungsanalyse durchgeführt, auf deren Grundlage verschiedene Konstruktions- und Antriebskonzepte bewertet wurden. Mithilfe einer Trade-off-Analyse fiel die Entscheidung auf einen dreirädrigen Roboter mit Differenzialantrieb, da dieses Konzept eine hohe Wendigkeit sowie einen vergleichsweise einfachen mechanischen Aufbau bietet. Gleichzeitig wurde eine modulare Systemarchitektur entwickelt, bei der die Bewegungssteuerung und die Benutzerinteraktion auf zwei getrennte Recheneinheiten aufgeteilt wurden. Während ein ESP32 die Motorsteuerung übernimmt, dient ein Raspberry Pi als Plattform für das Display und zukünftige Interaktionsfunktionen.  

Neben der elektrischen und softwareseitigen Konzeption wurde bereits ein erster mechanischer Prototyp gefertigt. Dieser bestand überwiegend aus einer Holzkonstruktion und diente der Erprobung des Fahrwerks sowie der Validierung der ausgewählten Komponenten. Darüber hinaus wurde die Motorsteuerung vollständig implementiert, sodass der Roboter über eine Fernsteuerung bewegt werden konnte. Die entwickelte Differenzialsteuerung ermöglichte sowohl Geradeausfahrten als auch Drehungen auf der Stelle. Zusätzlich wurden Sicherheitsmechanismen, wie ein Failsafe bei Verbindungsabbruch, implementiert und erfolgreich getestet.  

Ein weiterer Schwerpunkt der vorherigen Studienarbeit lag auf der Entwicklung der Mensch-Maschine-Interaktion. Hierfür wurde ein Display in das System integriert, auf dem animierte Augen sowie Statusinformationen dargestellt werden sollten. Die grundlegende Softwarearchitektur zur Ansteuerung des Displays sowie erste Animationen wurden erfolgreich umgesetzt und in Betrieb genommen. Die Integration weiterer Interaktionsmöglichkeiten sowie die endgültige Gestaltung des Roboters waren jedoch nicht Bestandteil der damaligen Arbeit.  

Zum Abschluss der ersten Projektphase lag somit ein funktionsfähiger Demonstrator vor, dessen Fahrwerk, Antriebssystem, Steuerung sowie grundlegende Interaktionsfunktionen erfolgreich validiert werden konnten. Die mechanische Konstruktion entsprach jedoch noch einem funktionalen Holzprototyp und war weder für den dauerhaften Einsatz noch hinsichtlich Design, Gewicht oder Modularität optimiert. Ebenso fehlten die vollständige Integration des Kopfes, ein finales Gehäusedesign sowie weitere Sensorik und autonome Funktionen. Diese offenen Punkte bildeten den Ausgangszustand und zugleich die Grundlage für die im Rahmen dieser Studienarbeit durchgeführte Weiterentwicklung des Roboters.  